%----------------------------------------------------------------------
%  Metadata — EDIT THIS FILE, nothing else in config/
%  thesis.tex passes these to the bossthesis class (cover page, title
%  page, declaration). scripts/build_review_docx.py reads the
%  \newcommand lines for the review copy.
%----------------------------------------------------------------------
\newcommand{\thesistitle}{Understanding Privacy Harm Communication in Health Tracking Apps}
\newcommand{\thesiskind}{Bachelor's Thesis in IT Security}
\newcommand{\thesiscitetype}{Bachelor's Thesis}
\newcommand{\authorname}{Hamza Sghaier}
\newcommand{\matriculation}{108022238911}
\newcommand{\birthdate}{\todo{date of birth}}
\newcommand{\birthplacename}{\todo{place of birth}}
\newcommand{\advisorname}{Dr. Kavous Salehzadeh Niksirat}
\newcommand{\examinerone}{Prof. Dr. Carmela Troncoso\todo{proposed in the registration form; confirm}}
\newcommand{\examinertwo}{\todo{second examiner}}
\newcommand{\thechair}{Max Planck Institute for Security and Privacy (MPI-SP)}
\newcommand{\thechairaddress}{Max Planck Institute for Security and Privacy\\Universitätsstraße 140 $\cdot$ 44799 Bochum $\cdot$ Germany}
\newcommand{\thefaculty}{Faculty of Computer Science}
\newcommand{\theuniversity}{Ruhr-Universität Bochum}
\newcommand{\theprogram}{B.Sc. IT-Sicherheit}
\newcommand{\startdate}{\todo{start date set by the examination office}}
\newcommand{\submissiondate}{\todo{submission date}}

% AI tools named in the declaration box. Drafted from what the repository
% records; the author must check, correct and complete every line.
\newcommand{\aitoollist}{%
\begin{itemize}
  \item \textit{Claude (Anthropic), used through Claude Code}: writing and
    debugging the collection, ranking and checking scripts in
    \texttt{scripts/} and the generated tables; preparing the eligibility
    screening decisions of the app sample, which the author checks
    (Section~\ref{sec:results-sample}); revising and drafting passages of
    Chapters~4 and~5 and Appendices~C and~E. \todo{Check and complete; name
    the sections.}
  \item \textit{OpenAI Codex}: \todo{the repository contains a Codex
    configuration; state whether and for what it was used.}
  \item \textit{Methods digest of the literature}: \todo{name the tool that
    produced the cards and summaries in 01\_Literature/AI Literature
    Analysis, and its use.}
  \item \textit{Language model of the coding pipeline}: part of the method,
    described in Chapter~\ref{ch:methodology} and Appendix~\ref{app:prompts}.
\end{itemize}}
